% Phase 6 (Andrés): Part 2 e-i. Every number here comes from src/thermodynamics.py (python -m src.thermodynamics).
\subsection{Entropy and thermodynamic implications}
With Boltzmann's entropy $S=k_B\ln\Omega$, the entropy difference between two macrostates of the same $N$ is
\begin{gather}
\Delta S_{k_1k_2}\equiv S(k_1)-S(k_2)=k_B\,[\ln\Omega(N,k_1)-\ln\Omega(N,k_2)],\\
\frac{P(k_1)}{P(k_2)}=e^{\Delta S_{k_1k_2}/k_B},
\end{gather}
computed with \texttt{delta\_entropy} (Phase 4) on top of \texttt{ln\_omega} (Phase 3), and reported in units of $k_B$ and in J/K with $k_B=1.380649\times10^{-23}$\,J/K (exact in the SI). The multiplicities are exact integers, so the only numerical error is floating-point rounding, below $10^{-12}$ in relative terms. Since $N=121$ is odd, the most probable macrostate $N/2=60.5$ is the degenerate pair $k=60,61$ with the same $\Omega$; we read ``$N/2+d$'' as $k=60+d$ and also give the other neighbour, $k=61+d$, which brackets the half-integer value.

\begin{table}[htbp]\centering\small
\caption{Entropy differences for $N=121$ coins. The last column is how many times more likely $k_1$ is than $k_2$.}\label{tab:entropy}
\begin{tabular}{@{}lrrrrr@{}}
\toprule
 & $k_1$ & $k_2$ & $\Delta S/k_B$ & $\Delta S$ (J/K) & $P(k_1)/P(k_2)$ \\
\midrule
(e) $N/2+25$ vs $N/2$ & 85 & 60 & $-10.1191$ & $-1.3971\times10^{-22}$ & $4.03\times10^{-5}$ \\
 & 86 & 60 & $-10.9899$ & $-1.5173\times10^{-22}$ & $1.69\times10^{-5}$ \\
(f) $N/2+10$ vs $N/2$ & 70 & 60 & $-1.4814$ & $-2.0453\times10^{-23}$ & $0.227$ \\
 & 71 & 60 & $-1.8123$ & $-2.5021\times10^{-23}$ & $0.163$ \\
(g) $N/2$ vs $20$ & 60 & 20 & $+29.3229$ & $+4.0485\times10^{-22}$ & $5.43\times10^{12}$ \\
\bottomrule
\end{tabular}
\end{table}

\paragraph{(e), (f), (g).} Table~\ref{tab:entropy}. Moving $25$ heads away from the peak lowers the entropy by $\approx10$--$11\,k_B$, so that macrostate is $\sim10^{-5}$ times as likely as the peak; $10$ heads away costs only $1.5$--$1.8\,k_B$, a factor $0.16$--$0.23$, so such fluctuations are common at $N=121$. In (g) the sign flips: going from the ordered state $k=20$ to the peak raises the entropy by $29.3\,k_B$, making the peak $5.4\times10^{12}$ times more likely. These values follow the Gaussian estimate $\Delta S/k_B\simeq-2[(k_1-N/2)^2-(k_2-N/2)^2]/N$ near the peak ($-9.92$ for (e), $-1.49$ for (f)) and depart from it further out ($27.1$ against the exact $29.3$ for (g)), where the distribution is no longer Gaussian. In joules per kelvin all these differences are of order $10^{-22}$\,J/K, tiny compared with macroscopic entropies because $N$ is only $121$.

\paragraph{(h) Waiting for 8 heads.} With a rate of $N$ coins per second, i.e.\ one toss of all $121$ coins every second, each toss lands on $k=8$ with $P(121,8)=3.384481\times10^{-25}$. The tosses are independent, so the number of tosses until the first success is geometric with mean $1/P$:
\begin{equation}
\langle t\rangle=\frac{1}{P(121,8)}=2.9547\times10^{24}\ \text{s}=(6.791\pm0.010)\times10^{6}\ \text{ages of the Universe},
\end{equation}
taking the age of the Universe as $13.787\pm0.020$\,Gyr $=4.351\times10^{17}$\,s (Planck 2018); its uncertainty is the only one in the result. Reading ``only 8 heads'' as ``at most 8 heads'' hardly changes it: $P(k\le8)=3.637227\times10^{-25}$ and $\langle t\rangle=2.749\times10^{24}$\,s $\approx6.3\times10^{6}$ ages of the Universe. Such an ordered state is allowed but never observed in practice.

\paragraph{(i) Most probable macrostate against $M\pm10$ heads.} Read as written, with $M=21$, the comparison is with $k=11$ and $k=31$:
\begin{equation}
\frac{P(60)}{P(11)}=1.501045\times10^{20},\qquad \frac{P(60)}{P(31)}=2.892531\times10^{6},
\end{equation}
($P(60)=0.0720892$, $P(11)=4.80\times10^{-22}$, $P(31)=2.49\times10^{-8}$). The two ratios differ by fourteen orders of magnitude because $k=11$ and $k=31$ lie $49.5$ and $29.5$ heads from the peak, and the suppression grows like the square of that distance. If the intended reading is $N/2\pm10$ instead, the comparison is with the mirror pair $k=51$ and $k=70$ ($51+70=N$, same $\Omega$), and the peak is only $4.399168$ times more likely. The conclusion is the same in both readings, only its strength changes: the peak is always the most probable macrostate, overwhelmingly so against states far from $N/2$ and only by a small factor against states within $\sim2\sigma$ ($\sigma=\sqrt N/2=5.5$).
